\documentclass[reprint,amsmath,amssymb,aps,prl,nofootinbib,floatfix]{revtex4-2}

\usepackage{graphicx}
\usepackage{bm}
\usepackage{physics}
\usepackage{placeins}
\usepackage[nopatch=footnote]{microtype}

% Reduce minor underfull/overfull boxes without manual line-breaking
\setlength{\emergencystretch}{2em}

\begin{document}

\title{Companion Draft: Emergent Spacetime Kinematics from Irreversible Memory Dynamics}

\author{H. Horn}
\affiliation{Peine Stederdorf, Germany}

\date{\today}

\begin{abstract}
This companion draft records a future-work programme for testing whether a minimal irreversible, finite-memory dynamics can support effective space and relativistic kinematics as coarse-grained consistency conditions of metastable, localized structures (``dynamical knots''). It does not claim a derivation of spatial dimensionality, finite signal speed, or Lorentz invariance. Those statements remain conditional targets for later multi-knot propagation and operator-based metastability tests.
\end{abstract}

\maketitle


%-------------------------------------------------
\section{Introduction}\label{sec:introduction}
%-------------------------------------------------

Modern physics treats space, time, and relativistic structure as fundamental ingredients. Even in approaches that aim at deeper unification--ranging from quantum gravity programs \cite{Rovelli2004Quantum} to discrete or information-theoretic models \cite{Bombelli1987Causal,Wen2004Quantum}--the dimensionality of space and the kinematical structure of spacetime are typically assumed from the outset rather than derived.

In Paper I, we introduced a minimal irreversible dynamical framework based on a single generative process with finite, relaxing memory \cite{Mori1965Transport,Zwanzig2001Nonequilibrium,Hanggi1995Colored}. Without assuming spacetime, geometry, or quantum postulates, we showed that an intrinsic arrow of update order arises from the non-invertibility of the memory update, that a coarse-grained notion of time emerges operationally, and that the dynamics self-organizes into localized, metastable structures (``dynamical knots'') characterized by intrinsic relaxation scales.

The present work builds directly on this foundation. Its purpose is to address the next conceptual step: under what conditions may the abstract state space $\Omega$ of the minimal model be interpreted as physical space \cite{Coifman2006Diffusion,Ollivier2009Ricci}, and why does a relativistic kinematics emerge at large scales?

Crucially, we do not assume spatial dimensionality, isotropy, or Lorentz invariance \cite{Liberati2013Lorentz}. Instead, we treat these notions as diagnostic properties of metastable dynamical regimes. Our guiding principle is that physical structure should be inferred from the internal consistency and robustness of the emergent dynamics, not imposed as a background scaffold.

The present work adopts a deliberately minimal and dynamical notion of structure. No geometry and no spacetime symmetry algebra are postulated. The only algebraic layer used from the start is the one already forced by Paper~I: the augmented Markov state and its one-step action on observables. In particular, spatial dimension, isotropy, and relativistic kinematics are not assumed but diagnosed operationally from the stability properties of metastable dynamical regimes. This places our approach conceptually close to a broad class of pre-geometric and emergent-geometry programs, while remaining distinct in its methodological commitments.

From a mathematical perspective, it is worth noting that modern geometry itself has increasingly moved away from treating classical manifolds as fundamental. Examples include noncommutative geometry, topos-theoretic approaches, and more recently perfectoid and adic spaces \cite{Scholze2012Perfectoid}, which are designed to interpolate between discrete and continuous structures across scales. While such frameworks demonstrate that classical geometric intuition need not be fundamental, we do not employ any of these formalisms directly. Our construction remains intentionally pre-geometric and dynamical, relying only on stochastic evolution with finite, relaxing memory.

The role of these developments in the present context is therefore purely conceptual. They provide independent evidence that emergent smooth geometry from non-classical substrates is neither exotic nor inconsistent. However, rather than introducing sophisticated mathematical structures at the outset, we ask which geometric notions are forced upon us by the internal consistency of the dynamics itself. In this sense, geometry is treated not as an axiom, but as an effective descriptor whose validity is limited to metastable, macroscopic regimes.

The working conjecture of this companion draft is that space and relativistic kinematics may arise as consistency conditions of interacting, memory-driven dynamical structures. In particular, the following remain test targets rather than established results:
\begin{itemize}
\item Effective spatial dimensionality is tied to the stability and decoupling of degrees of freedom of metastable knots \cite{Ambjorn2005Spectral,Benedetti2009Spectral,Carlip2017Dimension}.
\item Finite memory and delayed self-interaction may support finite response speeds only under additional locality and multi-knot transfer assumptions.
\item Requiring consistency of descriptions across different localized dynamical regimes could motivate Lorentz-like kinematical invariance at large scales \cite{Volovik2003Universe,Liberati2013Lorentz}.
\end{itemize}

Importantly, the scope of this paper is deliberately limited. We address only kinematical aspects of spacetime emergence. Quantum dynamics, spinor structure, and gauge interactions are deferred to a companion work. Likewise, we do not attempt to construct a dynamical theory of gravity; curvature and spacetime dynamics are explicitly outside the present scope.

\paragraph{Paper outline.}
We first state how an effective spatial interpretation of $\Omega$ could be tested, then outline finite-propagation diagnostics for multi-knot response, and finally describe the consistency conditions that would be required before Lorentz-like kinematics could be claimed.


\section{From Abstract State Space to Effective Space}\label{sec:effective_space}
\subsection{Status of the state space $\Omega$}

In Paper~I the dynamics was formulated on an abstract continuous state space $\Omega$, equipped with minimal coordinate structure sufficient to define local coarse-graining and gradients.
No physical interpretation of $\Omega$ as space was assumed.
In particular, neither spatial dimensionality nor isotropy were postulated.

The aim of the present section is to clarify under which conditions $\Omega$ admits an \emph{effective spatial interpretation} in the coarse-grained, metastable regime of the dynamics. 
Throughout, ``space'' will be understood operationally, as a set of coordinates in which localized dynamical structures can be consistently compared.

\subsection{Memory-driven dynamics}
We briefly recall the minimal dynamical rules introduced in Paper~I.
The state at update step $n$ is given by $z_n = (x_n,\rho_n)$, where
$x_n \in \Omega$ denotes the position of the process and $\rho_n$ is a
finite, relaxing memory distribution.

The memory update rule reads
\begin{equation}
\rho_{n+1}(x)
= (1-\alpha)\,\rho_n(x)
+ \alpha\,G_\sigma(x-x_{n+1}),  \label{eq:memory_update}
\end{equation}
where $0<\alpha<1$ is the memory relaxation rate and $G_\sigma$ a localized kernel of width $\sigma$.
This update is explicitly non-invertible and introduces an intrinsic arrow of update order.

As in Paper~I, the visible coordinate sequence $\{x_n\}$ is generally non-Markovian, while the augmented sequence $\{z_n\}$ is Markov by construction.
Let $\Sigma$ denote the augmented state space and write the one-step transition kernel as
\begin{equation}
P(z,A):=\mathbb{P}\{z_{n+1}\in A\mid z_n=z\},
\qquad A\subseteq\Sigma .
\end{equation}
On bounded observables $f:\Sigma\to\mathbb{R}$ this defines the positive unital Markov/Koopman operator
\begin{equation}
(Uf)(z)=\int_\Sigma f(z')P(z,dz').
\end{equation}
Its iterates form a forward semigroup rather than, in general, a reversible group action or deterministic algebra automorphism.
This operator layer is not an additional spacetime axiom; it is the bookkeeping structure of the finite-memory dynamics and provides the global language in which metastable knots can be tested as slow modes or almost-invariant sets \cite{NoeNuske2013,FroylandLloydSantitissadeekorn2010,WilliamsKevrekidisRowley2015,KlusEtAl2018,Parzygnat2017,Fritz2020}.

On coarse-grained scales $t=\alpha n$, the trajectory dynamics may be written in effective form as \cite{Oksendal2003}
\begin{equation}
dx_t
= -\gamma\,\nabla\Phi[\rho_t](x_t)\,dt
+ \sqrt{2D}\,dW_t, \qquad \gamma:=\frac{\eta}{\alpha}, \label{eq:effective_SDE}
\end{equation}
with diffusion constant $D\sim\varepsilon^2/(2\alpha)$.

\subsection{Dynamical knots and local linearization}
In metastable regimes the dynamics self-organizes into localized regions of repeated visitation, termed \emph{dynamical knots}. 
Operationally, a knot corresponds to a region in which the memory density remains elevated over many persistence times.

Near such a metastable configuration $x_\ast$, the effective dynamics Eq.~\eqref{eq:effective_SDE} may be linearized, yielding an Ornstein--Uhlenbeck (OU)-type process \cite{OrnsteinUhlenbeck1930}
\begin{equation}
dx_t
= -\gamma\,H\,(x_t-x_\ast)\,dt + \sqrt{2D}\,dW_t, \label{eq:OU_local}
\end{equation}
where $H=\nabla\nabla\Phi|_{x_\ast}$ denotes the Hessian of the memory-induced potential.
The eigenvalues of $H$ define characteristic relaxation rates of perturbations around the knot.
This Hessian spectrum is a local diagnostic.
At finite lag, the corresponding global diagnostic is the spectrum of the transfer/Koopman operator restricted to an appropriate coarse-grained representation of $z_n$.
A robust knot should therefore be visible both locally, through slow relaxation modes of $H$, and globally, through a persistent slow spectral subspace or metastable membership of the Markov operator.

\subsection{Classification of dynamical knot types}

Metastable dynamical knots are not uniform objects.
Depending on the balance between stochastic exploration, memory accumulation,
and relaxation, qualitatively distinct knot types arise.
We classify knots according to their stability properties and dynamical role.

\paragraph{Definition (Local relaxation spectrum).}
Let $H=\nabla\nabla\Phi|_{x_\ast}$ denote the Hessian of the effective potential
at a metastable configuration $x_\ast$.
We write its ordered eigenvalues as
\begin{equation}
\lambda_1 \le \lambda_2 \le \dots \le \lambda_d .
\end{equation}
The associated relaxation times are
\begin{equation}
\tau_i \sim (\gamma \lambda_i)^{-1}.
\end{equation}

\paragraph{Type I: Transient knots.}
A knot is called transient if
\begin{equation}
\tau_i \lesssim \tau_{\mathrm{mem}}
\quad\text{for all } i.
\end{equation}

Such structures decay on timescales comparable to the memory persistence time and do not define stable effective coordinates.
They dominate microscopic and short-time dynamics but do not contribute to macroscopic geometry.

\paragraph{Type II: Marginal knots.}
A knot is marginal if there exists a subset of modes $\mathcal{I}$ such that
\begin{equation}
\tau_i \gg \tau_{\mathrm{mem}}
\quad\text{for } i\in\mathcal{I},
\qquad
\tau_j \lesssim \tau_{\mathrm{mem}}
\quad\text{for } j\notin\mathcal{I}.
\end{equation}

Marginal knots exhibit partial stability and may support transient effective coordinates.
They are sensitive to fluctuations and typically appear near dynamical phase boundaries.

\paragraph{Type III: Robust macroscopic knots.}
A knot is robust if
\begin{equation}
\tau_i \gg \tau_{\mathrm{mem}}
\quad\text{for exactly } i=1,\dots,d_{\mathrm{eff}},
\qquad
d_{\mathrm{eff}}\ge 1,
\end{equation}

and the remaining modes relax rapidly.
Only such knots admit long-lived, coarse-grained descriptions and define effective spatial coordinates.
In the following, we argue that robust macroscopic knots require $d_{\mathrm{eff}}=3$.
Equivalently, in the operator language these are the regimes expected to generate stable almost-invariant components or slow spectral modes at lags long compared with $\tau_{\mathrm{mem}}$ but short compared with macroscopic relaxation.

\paragraph{Type IV: Delocalized regimes.}
If no subset of modes satisfies $\tau_i\gg\tau_{\mathrm{mem}}$, the dynamics does not localize.
Such regimes correspond to purely exploratory motion and do not support macroscopic observables.

\paragraph{Remark (Hierarchy of relevance).}
Only Type~III knots contribute to emergent spatial structure, finite signal propagation, and Lorentz kinematics.
Type~I and Type~II knots play a role in microscopic dynamics and transitions, but decouple from macroscopic physics.

\subsection{Effective dimensionality}
The local relaxation spectrum of Eq.~\eqref{eq:OU_local} provides a natural diagnostic of effective dimensionality.
Directions associated with large eigenvalues of $H$ relax rapidly and decouple from long-time dynamics, while directions with comparable small eigenvalues remain dynamically relevant.

We therefore define the \emph{effective spatial dimension} of a knot as the number of weakly coupled modes with relaxation times exceeding the memory persistence scale.
This definition is inherently scale-dependent and allows for transient or microscopic excursions into higher-dimensional behavior, while favoring low-dimensional stable regimes at macroscopic scales.

\subsection{Emergence of isotropy}
Isotropy is not assumed at the level of the state space $\Omega$.
However, metastable knots that are robust under stochastic perturbations tend to suppress strong directional asymmetries. Anisotropic relaxation spectra are dynamically unstable, as noise-driven drift leads to reorganization of the knot configuration.

As a consequence, long-lived macroscopic regimes typically exhibit approximate isotropy in their effective coordinates.
Numerical evidence indicates that three effective spatial dimensions form a particularly robust macroscopic regime, whereas higher-dimensional structures tend to be short-lived and lower-dimensional ones lack sufficient degrees of freedom for stable interactions.


\subsection{Why three dimensions (macroscopically)}

\paragraph{Status note.}
The current repository does not yet establish a robust dimension-selection theorem. The historical occupancy-dimension evidence and the current scalar memory-cloud result are useful diagnostics, but the claim of macroscopic $d=3$ selection remains future work.

The linearized OU description introduced above provides a diagnostic of stability for a given number of effective degrees of freedom.
It does not, however, by itself select a preferred spatial dimension.
A possible macroscopic selection of three dimensions would have to arise at an earlier, nonlinear stage of the memory-driven dynamics.

We consider a stochastic trajectory $\{x_n\}_{n\in\mathbb{N}}$ in an abstract state space $\Omega\simeq\mathbb{R}^d$ with finite-memory self-interaction.
The position update is
\begin{equation}
x_{n+1} = x_n + \varepsilon\,\xi_n - \eta\,\nabla\Phi_n(x_n),
\label{eq:dyn}
\end{equation}
with memory update
\begin{equation}
\rho_{n+1} = (1-\alpha)\rho_n + \alpha\,G_\sigma(\cdot-x_{n+1}),
\label{eq:mem}
\end{equation}
where $G_\sigma$ is localized and normalized.
No additional geometric structure is assumed beyond measurability.

\paragraph{Definition (Metastable knot).}
A region $K\subset\Omega$ is called metastable if
\begin{equation}
\liminf_{N\to\infty}\frac{1}{N}\sum_{n=1}^N\mathbf{1}_K(x_n) > 0,
\label{eq:knot_def}
\end{equation}
and the associated memory-induced drift balances stochastic fluctuations.

\paragraph{Lemma 1 (Recurrence obstruction for $d\leq 2$).}
Let $\{x_n\}$ be the stochastic component of \eqref{eq:dyn}.
For $d\leq 2$ the process is recurrent \cite{Polya1921}:
\begin{equation}
\mathbb{P}\bigl(x_n\in B_r(x_0)\ \text{i.o.}\bigr)=1
\qquad\forall r>0.
\label{eq:recurrence}
\end{equation}

\paragraph{Consequence.}
For $d\leq 2$ the accumulated memory gradient
$\nabla\Phi_n(x)$ diverges along recurrent returns,
\begin{equation}
\|\nabla\Phi_n\|\xrightarrow[n\to\infty]{}\infty,
\end{equation}
destroying the balance condition required by \eqref{eq:knot_def}.
Hence, no metastable knots exist for $d\leq 2$.

\paragraph{Lemma 2 (Memory dilution for large $d$).}
Let $B_\sigma(x)$ denote the support of $G_\sigma$.
Its volume scales as
\begin{equation}
\mathrm{Vol}(B_\sigma)\sim\sigma^d.
\label{eq:volume}
\end{equation}

The typical memory density obeys
$\rho(x)\sim\frac{N_{\mathrm{visits}}(x)}{\sigma^d}$, implying for the drift magnitude
\begin{equation}
\|\nabla\Phi\| \sim \|\nabla K\|\,\rho \propto \sigma^{1-d}.
\label{eq:dilution}
\end{equation}

\paragraph{Lemma 3 (Loss of confinement for $d\geq 4$).}
The stochastic term in \eqref{eq:dyn} induces fluctuations of order $\delta x_{\mathrm{noise}}\sim\varepsilon$.

Metastability requires
\begin{equation}
\|\nabla\Phi\| \gtrsim \varepsilon,
\label{eq:stability_cond}
\end{equation}
which fails for sufficiently large $d$ due to \eqref{eq:dilution}.
In particular, for $d\geq 4$ the inequality \eqref{eq:stability_cond} cannot be satisfied at fixed $(\sigma,\varepsilon)$.
Thus, no long-lived knots exist for $d\geq 4$.

Metastable macroscopic structure requires simultaneously
\begin{align}
d &\geq 3 \quad\text{(to avoid recurrent self-blocking)},
\qquad \\
&\leq 3 \quad\text{(to prevent memory dilution)}.
\end{align}
The unique solution satisfying both requirements is
\begin{equation}
\boxed{d = 3}.
\end{equation}

\paragraph{Remark (Scale dependence).}
Additional weakly coupled modes may exist on microscopic timescales.
Such modes satisfy
\begin{equation}
\tau_{\mathrm{mode}}\ll\tau_{\mathrm{mem}}\sim\alpha^{-1},
\end{equation}
and therefore do not contribute to macroscopic metastability.
The effective dimension is thus scale-dependent, with $d=3$ selected in the long-lived regime.

\subsection{Interim conclusion}
Physical space is thus identified not with the abstract state space $\Omega$ itself, but with the effective coordinate system defined by stable, approximately isotropic, low-dimensional metastable dynamical regimes.
Spatial dimensionality emerges as a diagnostic property of the dynamics rather than as a fundamental assumption.


\section{Finite Propagation Speed from Memory Coupling}\label{sec:finite_speed}

\subsection{Motivation and scope}
In Section~\ref{sec:effective_space} we argued that stable, isotropic, low-dimensional metastable regimes motivate an effective spatial interpretation of the abstract state space $\Omega$.
We now turn to the question of signal propagation: how do distinct localized dynamical structures influence one another, and at what rate?

Importantly, we do not assume locality, causality, or a light-cone structure a priori.
Instead, we analyze how information transfer arises when multiple memory-driven processes are coupled only through finite, relaxing memory fields.

\subsection{Coupled dynamical knots}

We now consider a collection of metastable dynamical knots.
Each knot is localized in the effective space identified in Section~\ref{sec:effective_space} and is characterized by its internal relaxation spectrum.

Interactions between knots arise exclusively through the memory field.
No direct instantaneous coupling between positions is assumed; all influence is mediated by modifications of the memory distribution.

\subsection{Retarded response from finite memory}
Because the memory update Eq.~\eqref{eq:memory_update} is finite and relaxing, the response of one knot to perturbations of another cannot be instantaneous.
At the coarse-grained level, variations induced by knot $A$ enter the effective potential experienced by knot $B$ through a retarded memory kernel \cite{Kubo1957},
\begin{equation}
\delta\Phi_B(t)
\sim \int_{-\infty}^{t} K_{AB}(t-t')\,\delta x_A(t')\,dt',
\label{eq:retarded_kernel}
\end{equation}
where $K_{AB}$ decays on a characteristic memory time scale.
This form expresses the fact that information transfer is delayed by the finite persistence of memory.

\subsection{Inconsistency of instantaneous influence}

If instantaneous propagation were allowed, Eq.~\eqref{eq:retarded_kernel} would reduce to a delta-like kernel.
Such behavior would contradict the finite relaxation time encoded in Eq.~\eqref{eq:memory_update}, effectively bypassing the irreversible memory dynamics that defines the model.

Consistency of the coupled dynamics would exclude instantaneous influence only under explicit locality assumptions and would then enforce a finite response time between spatially separated knots.

\subsection{Effective maximal propagation speed}
A delayed memory response can imply a maximal propagation speed for dynamical influence only when direct long-range coupling terms are absent or sufficiently suppressed.
On coarse-grained scales this speed may be estimated as
\begin{equation}
c_{\mathrm{eff}}
\sim \frac{\ell_{\mathrm{mem}}}{\tau_{\mathrm{mem}}},
\label{eq:ceff}
\end{equation}
where $\ell_{\mathrm{mem}}$ characterizes the spatial range of the memory kernel and $\tau_{\mathrm{mem}}\sim\alpha^{-1}$ its relaxation time.

The quantity $c_{\mathrm{eff}}$ is not a fundamental constant, but an emergent scale determined by the underlying memory dynamics.

\paragraph{Remark (Minimal assumptions behind a speed limit).}
The estimate \eqref{eq:ceff} becomes a genuine \,\emph{upper bound} once one assumes that (i) the memory kernel has a finite effective range $\ell_{\mathrm{mem}}$ (e.g. sufficiently fast decay), and (ii) influence between distant knots is mediated only through successive updates of the same relaxing memory field \cite{LiebRobinson1972}.
Under these conditions, a perturbation produced at knot $A$ cannot affect knot $B$ at distance $L$ in a time shorter than $\sim (L/\ell_{\mathrm{mem}})\,\tau_{\mathrm{mem}}$.
If, by contrast, the kernel had effectively infinite range or instantaneous nonlocal components, no strict finite-speed statement would follow.

\subsection{Universality}
For multiple knots sharing a common memory substrate, consistency of their effective descriptions would require any measured maximal propagation speed to be universal within the tested regime.
Different values of $c_{\mathrm{eff}}$ would prevent a coherent comparison of spatial and temporal scales between knots.

Thus, a universal maximal propagation speed remains a dynamical consistency target for interacting memory-driven structures.

\subsection{Interim conclusion}
Finite, relaxing memory alone does not generically enforce a finite and universal maximal propagation speed; this would require local transfer assumptions and multi-knot response evidence.
This is a conditional test target rather than a result: without local transfer assumptions and multi-knot evidence, finite memory alone does not prove a speed limit.

\section{Emergent Lorentz Kinematics}\label{sec:lorentz}
\subsection{Operational notion of observers}
In the metastable regime identified in Sections~\ref{sec:effective_space} and \ref{sec:finite_speed}, each dynamical knot defines an effective local coordinate system through its coarse-grained trajectory statistics.
Different knots correspond to distinct operational descriptions of the same underlying memory-driven dynamics.

An ``observer'' is therefore identified with a choice of metastable reference knot together with its associated effective spatial and temporal coordinates.
No preferred observer is assumed.

\subsection{Universality of the maximal propagation speed}

Section~\ref{sec:finite_speed} outlined conditions under which a finite maximal propagation speed $c_{\mathrm{eff}}$ could arise as a consistency condition of finite, relaxing memory.
For multiple knots to admit a coherent mutual description, this speed must be observer-independent.

Let $(t,\mathbf{x})$ and $(t',\mathbf{x}')$ denote effective coordinates associated with two distinct metastable knots.
Consistency of information transfer requires
\begin{equation}
\frac{\|\Delta\mathbf{x}\|}{\Delta t}
\leq
c_{\mathrm{eff}}
\quad\Longleftrightarrow\quad
\frac{\|\Delta\mathbf{x}'\|}{\Delta t'}
\leq
c_{\mathrm{eff}}.
\label{eq:ceff_invariance}
\end{equation}

\subsection{Kinematical invariance group}
We now characterize the class of transformations between effective coordinate
systems that preserve the bound \eqref{eq:ceff_invariance}.
Consider linear transformations of the form
\begin{equation}
\begin{pmatrix}
c_{\mathrm{eff}}\,t' \\
\mathbf{x}'
\end{pmatrix}
= \Lambda
\begin{pmatrix}
c_{\mathrm{eff}}\,t \\
\mathbf{x}
\end{pmatrix},
\label{eq:linear_transform}
\end{equation}
where $\Lambda$ is an invertible $(d+1)\times(d+1)$ matrix.


\paragraph{Remark (From invariant speed to an invariant interval).}
To pass from the observer-independent bound \eqref{eq:ceff_invariance} to a specific invariance group, we assume that in the macroscopic, metastable regime
(i) transformations between effective observers are well-approximated as linear,
(ii) the effective space is homogeneous and isotropic, and
(iii) the ``null'' propagation defined by $\|\Delta\mathbf{x}\|=c_{\mathrm{eff}}\Delta t$ is mapped into itself.
Under these standard kinematical assumptions, preserving the $c_{\mathrm{eff}}$-cone is equivalent (up to an overall scale choice) to preserving a quadratic form of signature $(1,d)$.

Requiring invariance of the quadratic form
\begin{equation}
s^2
:=
c_{\mathrm{eff}}^2 t^2 - \|\mathbf{x}\|^2
\label{eq:interval}
\end{equation}
yields the constraint
\begin{equation}
\Lambda^{\mathrm T}
\begin{pmatrix} 1 & 0 \\
0 & -\mathbb{I}_d \end{pmatrix}\Lambda 
= \begin{pmatrix} 1 & 0 \\
0 & -\mathbb{I}_d \end{pmatrix}.
\label{eq:lorentz_condition}
\end{equation}

\subsection{Emergent Lorentz group}

Equation~\eqref{eq:lorentz_condition} defines the Lorentz group
$\mathrm{O}(1,d)$ associated with the invariant speed $c_{\mathrm{eff}}$ \cite{Einstein1905,Rindler2006}.
Thus, the transformations relating different effective observers are elements of $\mathrm{O}(1,d)$, restricted to the proper, orthochronous component and assuming a consistent time orientation.
In this sense the Lorentz group is the effective stabilizer of the macroscopic propagation cone selected by the metastable operator dynamics, not a microscopic input symmetry of $\Omega$.

\subsection{Interpretation}

The emergence of Lorentz kinematics in the present framework does not rely on postulated spacetime structure, causal axioms, or relativistic principles.
Instead, it follows uniquely from:
\begin{enumerate}
\item the measured existence of a finite, universal maximal propagation speed $c_{\mathrm{eff}}$,
\item the requirement that effective descriptions associated with different metastable knots remain mutually consistent.
\end{enumerate}

Lorentz symmetry thus appears as a kinematical redundancy of the coarse-grained description, not as a fundamental microscopic symmetry.

\subsection{Limiting case: absence of proper time}

The invariant interval \eqref{eq:interval} implies that trajectories satisfying
\begin{equation}
c_{\mathrm{eff}}^2\,dt^2 = \|\mathrm d\mathbf{x}\|^2
\label{eq:null}
\end{equation}
have vanishing proper time.
In the present framework, such trajectories correspond to modes that do not form metastable, memory-supported structures and therefore carry no intrinsic relaxation scale.
This provides an operational interpretation of lightlike propagation without introducing additional postulates.

\subsection{Interim conclusion}
Lorentz kinematics would have to be tested as the effective linear invariance structure compatible with a measured finite, universal propagation speed and observer equivalence in the metastable regime of memory-driven dynamics.
No microscopic spacetime assumptions are required.

\section{Outlook}\label{sec:outlook}

This companion draft has outlined how effective space and relativistic kinematics could be tested as coarse-grained consistency conditions of an irreversible, finite-memory dynamics. Metastable localized structures would have to furnish operational coordinates, local memory-mediated transfer would have to produce a universal maximal propagation speed $c_{\mathrm{eff}}$, and observer equivalence in the metastable regime would then motivate Lorentz-like kinematics.

The present scope is purely kinematical. We do not construct a dynamical theory of gravity, and we do not introduce quantum postulates here. Quantum dynamics and the basis for the Standard Model are deferred to Paper~III.

\paragraph{Take-home and near-term targets.}
\begin{itemize}
\item \textbf{Minimal ingredients.} A finite memory persistence time and finite effective memory range are not by themselves sufficient; finite propagation requires local transfer assumptions and absence of direct long-range response terms.
\item \textbf{Where Lorentz symmetry can fail.} At finite coarse-graining and finite memory length, small departures from exact Lorentz invariance are expected; these deviations should scale with ratios such as $\tau_{\mathrm{mem}}/\tau_{\mathrm{macro}}$ and $\ell_{\mathrm{mem}}/L$.
\item \textbf{Testable numerical diagnostics.} In simulations one can measure (i) the retarded response kernel $K_{AB}$ between knots, (ii) the emergent speed estimate $c_{\mathrm{eff}}\sim \ell_{\mathrm{mem}}/\tau_{\mathrm{mem}}$, (iii) the local stability spectrum of knots, and (iv) transfer-operator spectra and metastable memberships on the augmented states $z_n$.
\item \textbf{Open dynamical questions.} A next step is to measure whether interactions between multiple knots define a stable $c_{\mathrm{eff}}$, constrain admissible effective field content in the metastable regime, and determine the stabilizer/equivariance groups of robust knot classes.
\end{itemize}

\bibliographystyle{unsrtnat}
\bibliography{references2}
\end{document}
